Per a obtenir un camp elèctric vertical aproximadament uniforme de 5\,000~N/C i dirigit cap amunt, disposem de dues plaques metàl·liques paral·leles separades 10,0~mm, a les quals apliquem una diferència de potencial.

\begin{apartat}{1}
Feu un esquema del muntatge en què indiqueu el signe de la càrrega de cada placa i representeu-hi les línies del camp elèctric. Calculeu la diferència de potencial entre les plaques i justifiqueu el signe del resultat.
\end{apartat}

\begin{apartat}{1}
Dues partícules de pols, de 0,50~$\mathrm{\mu g}$ de massa cadascuna, es troben entre les dues plaques. Una de les partícules (A) queda suspesa en equilibri i l'altra (B) es mou amb una acceleració de 14,7~$\mathrm{m/s^2}$ cap avall. Determineu la càrrega elèctrica de cada partícula. Considereu que entre les plaques no hi ha aire.
\end{apartat}

\blocetiquetat{Dada:}{$g = 9,81\ \mathrm{m\,s^{-2}}$}
